% 本文件由 LaTeX 排版 Agent 根据有效 translation.json 自动生成。
% author=Tertullian; work_id=0403; title=论贞女的蒙头

\chapter{论贞女的蒙头}

% structure: p0001 source=h1 target=omitted reason=page-title-already-rendered-by-parent

% structure: p0002 source=h2 target=section reason=relative-heading-rank; base=h2

\section{第一章：宁愿诉诸真理而非习俗，真理在其发展中渐进}

我已经承担了为我个人意见所特有的劳苦，现在也要用拉丁文证明：我们的贞女从过了年龄转折点之后，就应该蒙头；这一做法是真理所要求的，而真理是任何人都不能凭借先占权来强加的——无论时间的长度、人物的影响，还是地区的特权，都不能。因为习俗多半是从某种无知或单纯中发源的；然后它被接连不断地确认成惯例，并如此维持下去，与真理相悖。但我们的主基督自称是真理，而非习俗。如果基督是永恒的，且先于一切，那么真理同样也是永存和古老的。因此，那些把本质上古老的东西当作新奇事物的人，当自己省察。与其说是新奇，不如说是真理在揭露异端。凡与真理相抵触的，就是异端，即便它是古老的习俗。另一方面，如果有人对某事无知，这无知源于他自己的缺陷。此外，凡是无知之事，都应当像被人认可的事一样，被仔细地\Emph{探究}。事实上，信仰的准则全然是唯一的，不可动摇，也不可更改；即相信唯一全能的天主，宇宙的创造者，以及祂的圣子耶稣基督，由童贞玛利亚诞生，在本丢•彼拉多下被钉十字架，第三天从死者中复活，被接到天上，如今坐在圣父的右边，将要来临审判生者死者，并且（灵魂与）肉身的复活。这一信仰的法则恒常不变，而其他随后出现的纪律与行为准则，却容许改正的\Quote{新奇}；天主的恩宠，也就是说，运行并进展直到终结。因为这样一种假设算什么呢：既然魔鬼不断运作，日益增加邪恶的诡计，天主的工程却应当停止，或者停止进步？然而主派遣护慰者的理由正是：由于人性的平庸无法一次吸纳所有的事，纪律应当由主的代理人圣神，一点一点地指导、安排，并导向完美。祂说：\Quote{然而，}祂说，\Quote{我还有许多事要告诉你们，但你们现在还不能担负；当那一位真理之神来到时，祂会引领你们进入全部真理，并要把将来的事报告给你们。}（\BibleRef{若 16:12-13}）此外，祂在别处也预先宣告了祂的这一工作。那么，护慰者的管理职务是什么呢？不就是纪律的指导、圣经的启示、理智的改革、向\Quote{更美好的事}的进步吗？没有事物是没有成长阶段的：一切事物都等待它们的时机。简言之，训道者说：\Quote{凡事都有定时。}（\BibleRef{训 3:1}）请看，受造物本身是如何一点一点地走向结果实的。先是种子，从种子生出嫩芽，从嫩芽挣扎出灌木；此后枝干和叶子逐渐强壮，而我们称之为树的整体扩展；接着是胚芽的膨胀，从胚芽绽开花朵，从花朵打开果实；那果实本身起初是粗糙的、不成形的，一点一点地，沿着发展的正轨，被培养到风味的醇熟。同样，义德——因为义德与创造的天主是同一个——起初是初级的，具有对天主的自然敬畏；从那一阶段，它通过法律和先知前进到幼年；从那一阶段，它通过福音进入青年的热忱；如今，通过护慰者，它正在定居于成熟。祂将是基督之后唯一被称作并尊敬为主师的；因为祂不是凭自己说话，而是按基督的命令。祂是唯一的首席，因为唯有祂继承基督。那些接受了祂的人，把真理置于习俗之前。那些听见祂甚至到现在仍在预言，而非在古时的人，命令贞女完全蒙头。\par

% structure: p0004 source=h2 target=section reason=relative-heading-rank; base=h2

\section{第二章：在进一步讨论之前，先筛审习俗本身的问题}

但我暂时不把这一做法归因于真理。暂且就当作习俗：这样我就可以同样用习俗来反对习俗。\par

在整个希腊及其某些野蛮省份，大多数教会都让她们的贞女蒙头。在这片（非洲）天空下，也有地方实行这一做法；免得有人把这一习俗归因于希腊或野蛮的外邦人。但我提出那些由宗徒或宗徒弟子建立的教会作为（典范）；而且我想，早于某些（无名创始人）之前。因此，那些教会（与其它教会一样）拥有同样的习俗权威（可求助）；它们摆出\Quote{时代}和\Quote{导师}的对峙队列，比这些后来的（教会）更早更多。我们应当遵守什么？选择什么？我们不能轻蔑地拒绝一种无法谴责的习俗，因为它并不\Quote{陌生}，而且我们并不是在\Quote{陌生人}中间发现它，而是在与我们一起分享和平的律法和兄弟名号的人们中间。他们和我们有同一个信仰、同一天主、同一个基督、同一个希望、同一种圣洗圣事；我一次性说完：我们是一个教会。因此，凡属于我们弟兄的，都是我们的；只是，身体将我们分开。\par

仍然，在这里（正如在一切实践多样、怀疑和不确定的情况中通常发生的那样），应当加以考察，以看出两种如此不同的习俗中，哪一种更符合天主的纪律。当然，应当选择那种使童贞女蒙上头纱的习俗，因为她们只为天主所知；她们（除了光荣必须求自天主、而非来自人之外）甚至应为自己的特权而脸红。你赞美一个童贞女比责备她更使她脸红；因为罪恶的额头更为坚硬，从罪恶本身中学会无耻。因为那种在展示童贞女的同时却歪曲她们的习俗，永远不会得到任何人的认可，除非是那些与童贞女本身性格相似的人。这样的眼睛会希望一个童贞女被看见，正如那个希望被看见的童贞女一样。同类的眼睛相互渴求。看见和被看见属于同一种贪欲。一个贞洁的男子看见童贞女会脸红，正如一个贞洁的童贞女被男子看见会脸红一样。\par

% structure: p0008 source=h2 target=section reason=relative-heading-rank; base=h2

\section{第三章　习俗的逐步发展及其后果。对真理的热切呼吁}

但是，连那些最贞洁的教师也不曾选择在习俗之间进行考察。然而，直到最近，在\Emph{我们}中间，两种习俗都相对漠然地被接纳共融。这件事留给了选择，每个童贞女可以自己选择蒙头或裸露，正如她可以选择一样（她有同等的自由），就如同结婚一样，结婚本身既不被强制也不被禁止。真理曾满足于与习俗达成协议，以便在习俗的名义下，它甚至可以部分地享受自己。但当辨别的能力开始进步，以至于授予任何一种方式的许可成为显示更好部分的方式时，善事的大敌——尤其更是善制度的大敌——立刻开始了他自己的工作。属人的童贞女四处走动，与属天主的童贞女对立，颜面完全裸露，被激发出鲁莽的大胆；而那些有能力向丈夫索取一些东西的女人，更不用说那种要求（啊），即她们的对手——尤其是因为她们是唯独基督的\Quote{婢女}而更\Quote{自由}——可以被交给她们，她们展示了\Emph{童贞女}的假象。\Quote{我们感到惊愕，}他们说，\Quote{因为别人行事不同（于我们）；}并且他们宁愿被\Quote{惊愕}也不愿被激发（到谦逊）。如果我没有弄错，\Quote{惊愕}是一个榜样，不是好的，而是坏的，倾向于有罪的建造。好的事物只使邪恶的心惊愕。如果谦逊、腼腆、藐视荣耀、渴望唯独取悦天主是好的，那么被这些好事\Quote{惊愕}的女人，就让她们学会承认自己的邪恶。因为如果不节制的人也说自己被节制的人\Quote{惊愕}了呢？节操要被召回吗？并且，恐怕多婚者被\Quote{惊愕}，就要否定一夫一妻吗？为什么后者反而不能抱怨说，炫耀的童贞之轻率、无礼对\Emph{他们}是\Quote{惊愕}？那么，为了这些市场化的受造物，贞洁的童贞女们要被拖进教堂，在公共场合被认出时脸红，在被揭开面纱时颤抖，仿佛她们被邀请去遭受强暴？因为她们甚至同样不愿意遭受这个。每一个尊贵童贞女的公开暴露（对她来说）都是遭受强暴：然而肉体暴力的痛苦是较小的（邪恶），因为它来自自然的职责。但是当童贞女本身的精神因被剥夺遮盖物而受到侵犯时，她学会了失去她所保持的东西。啊，亵渎的手，竟有胆量扯掉奉献给天主的服装！如果迫害者知道这件（服装）是一个童贞女所选择的，他还能做出更坏的事吗？你裸露了一个少女的头，她立刻在自己眼中完全不再是童贞女；她经历了一种改变！因此，真理，起来吧；起来，仿佛从你的忍耐中迸发出来！我不希望你为任何\Emph{习俗}辩护；因为如今就连你曾享受自己自由的习俗也受到冲击！证明你就是那个遮盖童贞女者。亲自解释你自己的圣经，那是习俗所不理解的；因为，如果她理解了，她根本就不会存在。\par

% structure: p0010 source=h2 target=section reason=relative-heading-rank; base=h2

\section{第四章　论从\BibleRef{格前 11:5-16}引出的论据}

但是，既然连从圣经中提出与真理相反的论证也成为习俗，就立刻有人向我们提出事实说：\Quote{宗徒在论及头巾时没有提到童贞女，而只提到了\InnerQuote{妇女}；然而，如果他愿意童贞女也被遮盖，他会连同所提到的\InnerQuote{妇女}一起论到\InnerQuote{童贞女}；正如}（我们的对手说）\Quote{在他论述婚姻的那段经文中，他也同样论及\InnerQuote{童贞女}应遵守什么。}因此（有人主张）\Quote{她们不被包括在蒙头的律法中，因为这条律法中没有提到她们；更确切地说，这正是不\Emph{蒙头}的起源，因为不被\Emph{提到}的人也不被\Emph{命令}。}\par

但是我们也同样反驳这一论证。因为那位在其他地方知道如何分别提及每一性别——我的意思是\Emph{童贞女}和\Emph{妇女}，即\Emph{非童贞女}——为了区别起见；在这些（段落）中，他并没有提到\Emph{童贞女}，指出（通过不作区分）状况的共通性。否则他在这里也可以区分\Emph{童贞女}和\Emph{妇女}，正如在别处他所说：\Quote{\Emph{妇女}和\Emph{童贞女}是有分别的。}因此，那些被他沉默跳过而未加区分的人，他已包括在另一类中。\par

然而，因为在那情形下，\BibleQuote{\Emph{女人}和\Emph{童贞女}是分开的}，并不等于说这一区分也在当前情形的讨论中发挥其庇护作用，如某些人所认为的那样。因为许多在其他场合说的话，在那些它𠆢没有被说出的场合是没有分量的，除非话题与另一场合相同，以至于同一说法已足够。但先前\Quote{童贞女}和\Quote{女人}的情形与当前问题是广泛\Quote{分开}的。他说：\BibleQuote{\Emph{女人}和\Emph{童贞女}是分开的}。为什么？因为\BibleQuote{未出嫁的}，即童贞女，\BibleQuote{挂虑主的事，为叫身体和灵魂都能圣洁；但已出嫁的}，即非童贞女，\BibleQuote{挂虑怎样悦乐自己的丈夫}。这就是对那\Quote{分开}的解释，在本段经文中并无位置；此处所宣布的既非关于婚姻，也非关于女人和童贞女的心思念虑，而是关于蒙头的头巾。圣神愿意在这事上没有任何区别，便以\Quote{女人}这一个名称同样包括童贞女；祂没有特别提名童贞女，就没有把她与女人分开，而没有分开，就把她与那没有分开的女人结合在一起。\par

那么，现在使用基本词而同时让其他（从属分类）隐含在这个词中，这在无需逐一区分（整体的各个部分）的情况下，难道是一种\Quote{新奇}吗？自然，简洁的表达风格既令人愉悦又必不可少，因为冗长的言辞既使人厌倦又徒劳无益。因此，我们也满足于那些本身包含对特殊之理解的普遍词语。那么我们继续看这个词本身。表示\Emph{自然}区别的词是\Emph{女性}。对于这个自然词，\Emph{普遍}的词是\Emph{女人}。而普遍的词，再往下，\Emph{特殊}的词是\Quote{童贞女}、\Quote{妻子}、\Quote{寡妇}，或任何其他名称，甚至包括生命各个阶段的名称也都附加于此。因此，\Emph{特殊}的从属于\Emph{普遍}（因为普遍是优先的），\Emph{后继}的从属于\Emph{先行}的，\Emph{部分}的从属于\Emph{整体}：每一项都隐含在它从属的那个词本身之中，并且由于被包含在其中而被其所指。因此，当\Emph{身体}被命名时，无论是\Emph{手}、\Emph{脚}还是任何一个\Emph{肢体}，都不需要再单独指出。如果你说\Emph{宇宙}，其中就包括天及其中的一切——太阳、月亮、星座、星辰——以及地、海和构成元素列表的一切。当你命名由一切所构成的那一位时，你便命名了一切。同样，通过命名\Quote{女人}，他就命名了凡属于女人的一切。\par

% structure: p0015 source=h2 target=section reason=relative-heading-rank; base=h2

\section{第五章 论\Quote{女人}一词，尤其是与\Person{厄娃}的关系}

但既然他们以如此方式使用\Quote{女人}之名，以为它只适用于那与男子有过交合的女子，那么我们就要证明这个词的恰切性是属于性别本身，而非性别的某一等级；以便童贞女也（与其他人一样）能被普遍地包括在内。\par

当这种第二人类被天主造来作为男人的助手时，那\Quote{女性}立即被命名为\Quote{女人}；她那时仍是幸福的，仍配得乐园，仍是\Quote{童贞女}。（\Person{亚当}）说：\BibleQuote{她将被称为女人}。因此你们拥有了这个名字——我强调，不是已经\Quote{共同}属于\Quote{童贞女}，而是——\Quote{专属}于她；这名字从起初就被赋予了一个\Quote{童贞女}。但有些人巧妙地认为这是对未来而言的，说\Quote{她\Emph{将}被称为\Emph{女人}}，仿佛她注定要在她放弃童贞之后才如此称呼；因为他又加上：\BibleQuote{为此，人要离开父亲和母亲，与自己的\Emph{女人}结合，二人成为一体}。因此，让那些持此微妙见解的人先向我们证明，如果她具有未来参照地被称为\Quote{女人}，那么她同时的名号是什么。因为没有表示她当下品质的名字，她不可能存在。但这是什么假说呢？一个人，为着未来而被一个确定的名字称呼，而在当下却没有一个姓氏？\Person{亚当}给所有动物命名；没有一个基于未来状态，而是基于每一特定本性所服务的当下目的；（它们）被那从起初就显明倾向的名称所称呼。那么，她当时被叫做什么？为什么，每当她在\WorkTitle{圣经}中被提及，她在出嫁之前就拥有\Quote{女人}的称呼，而从未在被称作\Quote{童贞女}时是\Quote{童贞女}？\par

这名称在当时是她唯一的名号，那时（关于未来）尚无任何预言性的话。因为当圣经记载\Quote{二人都赤身露体，就是亚当和他的\Emph{女人}}时，这并不带有未来的意味，仿佛它说\Quote{他的\Emph{女人}}是作为\Quote{\Emph{妻子}}的预兆；而是因为他的\Emph{女人}尚未婚配，正如她是出自他的本体。他说：\Quote{这块骨，从我骨中，肉从我肉中，要称为\Emph{女人}。}\BibleRef{创 2:23}因此，正是出于自然默识，灵魂的实在神性已将我们称\Emph{妻子}为\Emph{女人}的习惯，不知不觉地引入常用口语中（正如它也将许多其他事物引入行为与言语中，我们将在别处证明它们源于圣经）；尽管我们在某些情况下也说得不尽恰当。因为连希腊人比我们更常用\Emph{女人}一词表示\Emph{妻子}，但仍有其他专门称呼\Emph{妻子}的名字。然而我宁可把这用法视为圣经的见证。因为当两人藉婚姻纽带成为一体时，那\Quote{肉中肉、骨中骨}者，便被称为那人的\Emph{女人}，因为她开始被算作他的本体，成了他的\Emph{妻子}。因此，\Emph{女人}并非天生是\Emph{妻子}的名称，而\Emph{妻子}因处境而成为\Emph{女人}的名称。总之，\Emph{女人}性可以不依\Emph{妻子}性而存在，但\Emph{妻子}性却不能脱离\Emph{女人}性，因为它甚至无法存在。因此，既确立了那新造女性的名称——即\Emph{女人}——并解释了她原先是什么，也就是为这名称作证，他就立刻转向预言的意义，说：\Quote{为此，人要离开父亲和母亲。}\BibleRef{创 2:24}这名称与预言如此分离，正如预言与那女子本人分离一样；当然，他所讲（预言）并非针对厄娃本身，而是指向那些未来的女性，他已在女性种族的母源中给她们命名。况且，亚当不会为厄娃的缘故离开\Quote{父亲和母亲}——他本无父母。因此，那预言性的话并不适用于厄娃，因为它也不适用于亚当。因为那是针对丈夫们的处境而预言的，他们注定要为\Emph{女人}的缘故离开父母；这不可能应验在厄娃身上，同样也不可能应验在亚当身上。\par

倘若情况如此，那么很明显，她并非因未来的（情况）而被加称\Emph{女人}，因为那未来（的情况）与她无关。\par

此外，还有一点：（亚当）本人也公布了名称的理由。因为他说了\Quote{她要称为\Emph{女人}}之后，又说\Quote{因为她是从男人身上取出来的}\BibleRef{创 2:23}——而男人本身当时也仍是童贞。但我们将在适当处也谈论\Emph{男人}的名称。因此，任何人都不应以预言性含义来诠释一个从其他意义推衍出来的名称；尤其因为她在何时确实获得了基于未来（情况）的名称——即在那里，她以\Emph{个人}的名字被称为\Quote{\Person{厄娃}}，因为\Emph{自然}的名字已在前。因为如果\Quote{\Person{厄娃}}意为\Quote{众生的母亲}\BibleRef{创 3:20}，看，她是从未来（情况）而得名的！看，她被预先宣告为\Emph{妻子}，而非\Emph{童贞}！这将是那将要结婚者的名字；因为母亲出自新娘。\par

因此，在此也表明，她当时被称为\Emph{女人}并非出于未来（的情况），而她在不久之后将获得那适合其未来状况的名称。\par

对这部分（问题）的回答已经足够。\par

% structure: p0023 source=h2 target=section reason=relative-heading-rank; base=h2

\section{第六章 论及玛利亚的平行情况}

现在让我们看看宗徒是否也遵循这名称的规范，与创世纪一致，将其归于\Emph{性别}；称\Emph{童贞}玛利亚为\Emph{女人}，正如创世纪称厄娃为\Emph{女人}。因为他在写给迦拉达人的信中说道：\Quote{天主派遣了祂的圣子，由\Emph{女人}所生}\BibleRef{迦 4:4}，当然，公认她是\Emph{童贞}，尽管\Person{厄比翁}反对（这教义）。我也承认天使加俾额尔曾被派遣到\Quote{一位\Emph{童贞}}那里\BibleRef{路 1:26}。但当他祝福她时，是\Quote{在\Emph{女人}中}，而非在\Emph{童贞}女中，将她排列：\Quote{在\Emph{女人}中你是蒙祝福的。}\BibleRef{路 1:28}天使也知道，即使\Emph{童贞}女也被称为\Emph{女人}。\par

但对这两个（论点），又有一人自以为作出了巧妙的回答；（其意以为）既然玛利亚是\Quote{已聘之女}，因此天使和宗徒才称她为\Emph{女人}；因为\Quote{已聘之女}在某种意义上就是\Quote{新娘}。然而，\Quote{在某种意义上}与\Quote{真理}之间毕竟大不相同，至少在此处是这样：因为别处，我们承认，必须如此持守。但现在，他们称玛利亚为\Emph{女人}，并非因她已出嫁，而是因她即使未被聘娶，也不失为一位\Emph{女性}；这是自始就受此称呼的：因为那必具有一种预判的力量，常规的类型正是从它衍生而来的。否则，就本节而言，若将玛利亚等同于\Quote{已聘之女}，以致称她为\Emph{女人}不是基于她是\Emph{女性}，而是基于她被许配给一个丈夫，那么立刻就会得出基督并非由\Emph{贞女}所生，因为（由）一位\Quote{已聘之女}（所生），而此人将因此不再是\Emph{贞女}。反之，若祂是由\Emph{贞女}所生——尽管也\Quote{已聘}，却仍童贞——那就承认，即使一位\Emph{贞女}，即使一位童贞者，也被称为\Emph{女人}。在此，无论如何，绝无预言性的说法，好似宗徒在说\Quote{由\Emph{女人}所生}时，是称一位\Emph{未来的女人}，即\Emph{新娘}。因为他不能指称一位后来的\Emph{女人}，基督并非要由她而生——即一位认识男人的女人；而是当时在场的那位，她是一位\Emph{贞女}，却仍被称为\Emph{女人}，这是由于这名称的恰当性——按照原初的规范，被主张（属于）一位\Emph{贞女}，因而也属于\Emph{女人}的普遍类别。\par

% structure: p0026 source=h2 target=section reason=relative-heading-rank; base=h2

\section{第七章 论宗徒为禁止女人蒙头所举的理由}

接下来我们应探讨宗徒教导女人应当蒙头的理由本身，看看这些同样的理由是否也适用于\Emph{贞女}；这样，通过发现有相同的理由需要蒙头，\Emph{贞女}与\Emph{非贞女}之间的名称共性也可以得到确立。\par

如果\BibleQuote{男人是女人的头}，当然也是贞女的头，因为结婚的女人是由贞女而来的；除非贞女是第三种族类，一种自有其头的怪物。如果\BibleQuote{女人若剃发或剪发，便是羞耻}，当然对贞女也是如此。（因此，让与天主打对台的世界注意，若它声称短发对贞女是优雅的，正如长发对男孩是优雅的一样。）那么，对她而言，剃发或剪发同样不体面，蒙头也就同样合宜。如果\BibleQuote{女人是男人的光荣}，那么贞女更是如此，因为她同时也是自身的光荣。如果\BibleQuote{女人是由男人而出}，且\BibleQuote{是为了男人}，那么亚当的那条肋骨起初就是贞女。如果\BibleQuote{女人在头上应有权柄}，那么贞女更是理当如此，因为这一断言的原因本质属于她。因为如果是\BibleQuote{为了天使}——即我们所读到的那些因贪恋女性而从天主和天上坠落的天使——谁能推测这些天使所渴望的是已经玷污的身体和人类情欲的残余，而不是更为贞女所点燃，因为她们的青春同样为人类情欲提供借口？因为圣经也这样暗示：\BibleQuote{当人在地上开始增多，也生了女儿；天主的儿子们看见人的女儿美丽，就随意挑选，娶来为妻。}因为这里\Emph{女人}的希腊名称似乎具有\Emph{妻子}的含义，因为提到了婚姻。那么，当它说\Emph{人的女儿}时，显然是指\Emph{贞女}，她们仍被视为属于\Emph{父母}——因为\Emph{已婚妇女}被称为其\Emph{丈夫}的——而它本可以说\Emph{人的妻子}；同样，它不称天使为奸夫，而是称为丈夫，因为他们娶了\Emph{未婚}的\Emph{人的女儿}，这些女儿上文说是\Emph{所生的}，从而也表明了她们的\Emph{童贞}：首先是\Emph{所生的}，但这里是与天使成婚。我不知道她们除了是\Emph{所生的}和后来结婚之外还有别的什么。如此危险的面容应当被遮掩，它竟将绊脚石抛到天上：这样，当它站在天主面前——在天主的审判台前被指控致使天使被逐出他们（原有的）领域——它也会在其他天使面前羞愧，并抑制那先前头上邪恶的自由——这自由如今连在人眼前也不应展示。但即使这些天使所渴望的是已经玷污的女性，那么\BibleQuote{为了天使}，贞女更应当蒙头，因为贞女更可能是天使犯罪的原因。此外，如果宗徒进一步引用\Quote{本性}的预判，即丰盈的头发是女人的荣耀，因为头发可作遮盖，那么这特别是贞女的荣耀；因为她们的装饰正在于此：将头发聚拢在头顶，以头发环绕完全覆盖头的堡垒。\par

% structure: p0029 source=h2 target=section reason=relative-heading-rank; base=h2

\section{第八章 反论}

这些（考虑）的反面情况，总之，造成了\Emph{男人}不可遮盖头部：即，因为他没有天生被赋予过多的头发；因为剃头或剪短对他来说并不是羞耻；因为天使并非因他而犯罪；因为他的头是\Person{基督}。因此，既然宗徒在论述\Emph{男人}和\Emph{女人}——为何后者应当蒙头，而前者却不——他何以对\Emph{童贞女}保持沉默，就显而易见了；也就是说，他容许\Emph{童贞女}被理解为包含在\Emph{女人}之内，基于同样的理由，他也避免提及\Emph{男孩}，认为\Emph{男孩}已包含在\Emph{男人}之内；他用\Emph{女人}和\Emph{男人}这两个专有名称涵盖了各自性别的全部类别。同样地，\Person{亚当}在尚纯洁时，在\WorkTitle{创世纪}中被称为\Emph{男人}：\BibleQuote{她将被称为，}他说，\BibleQuote{\Emph{女人}，因为她取自她的\Emph{男人}}。如此，\Person{亚当}在婚姻结合之前就是一个\Emph{男人}，同样\Person{厄娃}是一个\Emph{女人}。宗徒已相当明确地将他的论述应用于每一性别的普遍类别；并且简短而完整地，用如此恰当的定义，他说：\BibleQuote{每一个\Emph{女人}}。什么是\BibleQuote{每一个}？岂不是每一阶级、每一等级、每一状况、每一尊位、每一年龄？——如果（如事实那样）\BibleQuote{每一个}意味着总和与整体，并且在任何部分都无欠缺。但\Emph{童贞女}毕竟是\Emph{女人}的一部分。同样，关于\Emph{男人}不蒙头，他说\BibleQuote{每一个}。看哪两个不同的名称，\Emph{男人}和\Emph{女人}——各自带有\BibleQuote{每一个}；两条法律，彼此区分；一边是蒙头的法律，另一边是露头的法律。因此，如果说\BibleQuote{每一个\Emph{男人}}这句话表明\Emph{男人}这个名称甚至也适用于尚未成为\Emph{男人}的男性\Emph{少男}；再者，既然\Emph{名称}按照本性是共同的，那么，按照纪律，不蒙头的法律也同样适用于\Emph{男人}中的\Emph{童贞者}：为什么由此不能预先断定，既然\Emph{女人}被命名，每一个\Emph{女人}\Emph{童贞女}也同样被包含在共同\Emph{名称}中，从而也被包含在共同\Emph{法律}中呢？如果\Emph{童贞女}不是\Emph{女人}，那么\Emph{少男}也不是\Emph{男人}。如果因\Emph{童贞女}不是\Emph{女人}而她不蒙头，那么就因\Emph{少男}不是\Emph{男人}而让他蒙头。让\Emph{童贞}的共同身份分享同等的豁免。正如\Emph{童贞女}不被强迫蒙头，同样\Emph{男孩}也不要被命令不蒙头。为什么我们部分承认宗徒的定义是绝对的，关于\BibleQuote{每一个\Emph{男人}}，而不进一步探究他为何也没有提到\Emph{男孩}；但部分地却诡辩，尽管它同样绝对地关于\BibleQuote{每一个\Emph{女人}}？他表明在这点上曾有过一些争论；为了平息争论，他使用了全部简洁的语言：一方面不提名\Emph{童贞女}，以示对她的蒙头无可置疑；另一方面，却提名\BibleQuote{每一个\Emph{女人}}，而如果争论仅限于\Emph{童贞女}，他本会提名\Emph{童贞女}。\BibleQuote{如果有人好争论，}他说，\BibleQuote{我们并没有这种习惯，天主的教会也没有。}同样，格林多人自己也如此理解他。事实上，直到今日格林多人仍为他们的\Emph{童贞女}蒙头。宗徒所教导的，他们的门徒予以认可。\par

% structure: p0031 source=h2 target=section reason=relative-heading-rank; base=h2

\section{第九章 蒙头与为童贞女和一般妇女所遵守的其他纪律规则一致}

现在让我们来看，正如我们已经证明从本性和事物本身得出的论点同样适用于\Emph{童贞女}（也适用于其他\Emph{女性}），那么同样地，关于\Emph{女人}的教会纪律的规诫是否也关注\Emph{童贞女}。\par

女人在教会中不准说话；也不准她教导、施洗、献祭，或要求自己分得任何男人职分中的一份，更不用说（任何）司铎职务了。我们查考一下，这些事对\Emph{童贞女}是否合法。如果对\Emph{童贞女}\Emph{不}合法，但她与\Emph{女人}受制于同样条件，并与\Emph{女人}一同被分派了谦卑的必要，那么，哪一件事会对\Emph{她}合法，而对任何别的\Emph{女性}却不合法呢？如果有谁是\Emph{童贞女}，并已决心圣化自己的肉体，她因此获得了什么超越自身处境的特权呢？难道允许她免去蒙头，是为了使她进入教会时显眼可辨？为了让她在头顶的自由中显示圣洁的尊荣？若赐予她某种男人等级或职务的特权，本可给她更卓越的区分！我确实知道，在某个地方，一位不到二十岁的\Emph{童贞女}竟被列在\Emph{寡妇}的等级中！倘若主教必须给她某种救济，他当然可以用其他方式而不损害对纪律的尊重；免得这种奇迹——更确切说是怪物——在教会中被指指点点，即一位\Emph{童贞寡妇！}其更为奇特的是，她甚至作为\Emph{寡妇}也没有蒙头；两种方式都否认了自己：作为\Emph{童贞女}，因为她被算作\Emph{寡妇}；作为\Emph{寡妇}，因为她被称为\Emph{童贞女}。但是，允许她露头坐在那座位上的权威，也正是允许她以童贞女身份坐在那里的同一权威：那个座位，（除了\Quote{六十年}不单是\Quote{一夫之妻}（\Emph{女人}）——即\Emph{已婚妇女}——最终被选上，而且还要是\Quote{母亲}，是的，还要是\Quote{教养子女的}；）为的是，她们在一切情感中的经验训练，一方面使她们能够轻易地用劝言和安慰帮助别人，另一方面她们也走过了女性可以经受的全部考验之路。确实如此，根据她的地位，\Emph{童贞女}在公共荣誉方面没有任何许可。\par

% structure: p0034 source=h2 target=section reason=relative-heading-rank; base=h2

\section{若女童贞如此显眼，为何男童贞不可？}

同样地，基于任何区别也不允许。否则，就很不礼貌：女性处处服从男人，她们在额前戴着童贞的可敬标志，因此受到弟兄们的仰望、凝视和赞美，而如此多的男童贞、如此多的自愿阉人，却要秘密地携带自己的荣耀，没有任何标记使他们也同样显赫。因为他们也将不得不为自己要求一些区别——要么是加拉曼特人的羽毛，要么是野蛮人的头带，要么是雅典人的蝉饰，要么是日耳曼人的卷发，要么是不列颠人的刺青；或者采取相反的方式，让他们蒙着头潜伏在教会中。我们确信，如果圣神曾向女性作出此类让步，祂理当更愿意向男性作出；因为除了性别的权威，男性更应因贞洁本身而受到尊荣，这也更相宜。他们的性别对女性越是热切和炽热，那么这种更大热情的贞洁所包含的劳苦就越多；因此，如果童贞的炫耀是尊严，那么它就更配得一切炫耀。因为贞洁难道不超越童贞吗？无论是寡妇的贞洁，还是那些经同意已经弃绝了婚姻所包含的共同耻辱之人的贞洁？因为童贞的恒常是靠恩宠维持的，而贞洁的恒常是靠德行维持的。因为当你已经习惯于某种情欲时，克服情欲的斗争是巨大的；而有一种情欲，你从未享受过它的满足，你将轻易克服，因为你不曾有过享受之情欲这一对手。那么，天主岂会未能向男人作出比女人更多的让步吗？无论是基于更亲密的关系——作为\Quote{祂自己的肖像}，还是基于更艰苦的劳苦？但如果对男性什么也没有让步，那么对女性就更不用说了。\par

% structure: p0036 source=h2 target=section reason=relative-heading-rank; base=h2

\section{蒙头之规不适用于孩童}

但我们在上面为了后续讨论而暂时搁置的问题——以免破坏其连贯性——我们现在将用一个回答来了结。因为当我们针对宗徒的绝对定义——即\Quote{\Emph{每个女人}}必须被理解为（作为\Emph{女人}）甚至每个年龄——提出争论时，对方可能会反驳说，既然如此，童贞女就应从出生起，从她年龄的最初进入（时间之卷）时就当蒙头。\par

但事实并非如此；而是从她开始自觉、醒悟到自己的本性、从\Emph{童贞女}的（意识）中浮现、并体验到属于下一个年龄的那新奇（感觉）之时起。因为人类的始祖亚当和厄娃，在无理智之时，是\Quote{赤裸}的；但自从他们尝了\Quote{认识树}之后，首先感受到的不过是羞耻的缘由。于是他们各自用遮盖物标示对自己性别的认识。但即使她蒙头是\Quote{为了天使的缘故}，无疑蒙头之律法开始生效的年龄，正是\Quote{人的女儿}能够招致他人对她们身体的贪欲并经历婚姻的年龄。因为一个\Emph{童贞女}从她有可能\Emph{不}是童贞女之时起，就不再是\Emph{童贞女}。因此，在以色列，除非有血的证明证实其成熟，否则不得将女子交付给丈夫；故在此迹象之前，本性尚未成熟。所以，如果她在未成熟时是\Emph{童贞女}，那么当她被认为成熟时，她就不再是\Emph{童贞女}；并且作为\Emph{非童贞女}，现在她服从于律法，正如她服从于婚姻一样。至于\Emph{已订婚的女子}，确实有黎贝加为榜样：当她被带往——她自己仍未知——一位未知的未婚夫那里时，一得知她从远处看见的那人就是那男人，她没有等待握手、亲吻或问候；而是坦承自己的感受——即她已在精神上成婚——便在当时当地蒙上头帕，否认自己是\Emph{童贞女}。啊，已经属于基督训导的\Emph{女人}！因为她表明，婚姻如同淫乱一样，是通过目光和意念完成的；只不过有些人至今仍蒙着\Emph{黎贝加}的头。至于其余的（即那些\Emph{未}订婚的），让父母因经济拮据或过于审慎而造成的拖延来照管她们吧；让守贞的誓愿本身来照管她们吧。这种拖延丝毫不涉及一个已经在走自己指定道路、并向成熟付出其应付款项的年龄。另一位隐秘的母亲——大自然，和另一位隐藏的父亲——时间，已将他们的女儿嫁给了自己的律法。看啊，你那\Emph{童贞女儿}已经出嫁——她的灵魂因期待，她的肉身因转变——而你却在为她预备第二个丈夫！她的声音已变，四肢已发育完全，她的\Quote{羞耻}处处包裹着自己，月月献上贡物；而你还要否认她是\Emph{女人}，而你声称她在经历\Emph{女性}的体验？如果\Emph{男人}的接触使一个女人成为\Emph{女人}，那么除了实际经历婚姻之后，就不应有任何遮盖。不，甚至在外邦人中，（已订婚的女子）也是蒙着\Emph{面纱}被领到丈夫面前。但若她们是在\Emph{订婚}时蒙头，因为那时她们在身体和灵性上都已通过亲吻和握手与男子结合，通过这些方式她们首先在精神上开启了她们的端庄，通过双方良心的共同承诺，她们彼此交付了全部的羞怯；那么，时间岂不更会使她们蒙头呢？——没有时间，她们不能成为婚配者；而由于时间的紧迫，即使没有婚约，她们也不再是\Emph{童贞女}。时间甚至外邦人也遵守，以便服从自然律法，将各年龄应得的权利归还。因为他们从十二岁起就打发\Emph{女性}去处理事务，而\Emph{男性}则晚两年；将青春期定为在年龄上，而非在订婚或婚礼上。有人被称为\Quote{主妇}，尽管她是\Emph{童贞女}，以及\Quote{家长}，尽管是年轻人。而\Emph{我们}甚至连自然\Emph{律法}也不遵守；仿佛大自然的天主是另一个人，而非我们的天主！\par

% structure: p0039 source=h2 target=section reason=relative-heading-rank; base=h2

\section{第十二章 女性身份不言自明，不可仅靠露头来掩饰}

要认出\Emph{女人}，是的，要认出\Emph{已婚女人}，凭她身体和灵性的证据，这些证据她在良心和肉体上都体验到。这些是\Emph{自然}订婚和婚礼的早期石版。从外部给那个（已）有内部遮蔽的女子戴上头纱。让下身不裸露的女子上身也同样遮盖。你想知道年龄带有何种权威吗？将（这两种）各放在自己面前：一个过早地被挤压在\Emph{女性}的服装中，另一个虽然已成熟，却仍以适当的服装坚持\Emph{童贞}：前者更容易被否认是\Emph{女人}，而后者更难被相信是\Emph{童贞女}。那么，年龄的诚实就是这样，甚至连服装也无法压倒它。至于我们的这些（\Emph{童贞女}）甚至\Emph{凭借}服装就坦承她们年龄的变化；并且，一旦她们认识到自己是\Emph{女人}，就离开\Emph{童贞女}的行列，放下（从头开始）她们先前的自我：染发；用更放纵的发针固定头发；头发从前面分开，表明明显的\Emph{女性身份}。接下来，她们照镜子以助美貌，用洗脸来减弱过分的脸，也许还用化妆品来修饰，潇洒地披上斗篷，紧紧穿上各式鞋子，携带更充足的用品去浴场。我为什么要逐一列举呢？但她们明显的用品本身就展示了她们完全的\Emph{女性身份}：然而她们却想仅凭露头一事来扮演\Emph{童贞女}——用一个特征否认她们通过整个举止所表明的一切。\par

% structure: p0041 source=h2 target=section reason=relative-heading-rank; base=h2

\section{第十三章 若露头是适当的，为何不常在教会内外都这样做呢？}

若她们为了人的缘故而穿着虚假的服饰，那么即便为此目的，也当完全履行那服饰；既在教外人面前蒙头，就至少应在\Emph{教会内}隐藏她们的童贞，因为她们在教会外也是如此。她们害怕陌生人：也当敬畏弟兄们；否则，就该有前后一致的胆量，在街市上也如\Emph{童贞女}般出现，如同她们在教会中那样。若她们能在教外人那里卖掉一丝童贞，我倒要称赞她们的勇气。无论在异乡还是家中，身份相同；在主面前与在人面前，习俗相同，在于自由相同。那么，她们为何在异乡将荣耀藏匿，却在教会中暴露？我要问一个理由。是为悦乐弟兄们，还是天主自己？若为天主自己，祂既能看见暗中做的一切，也公义地赏报只为祂荣耀而行的事。最后，祂命令我们不可张扬任何一件在祂面前堪受赏报的事，也不可从人那里获取报酬。但如果我们被禁止让\Quote{左手知道}我们施舍一文钱或任何施舍之物，那么当我们向天主奉献如此伟大的祭品——我们的身体和灵魂——献上我们的本性时，我们该当隐藏在多深的黑暗中！因此，凡不能看似为天主而行的事（因天主不愿如此行），便是为人而行——这当然首要地是不法的，因为它暴露了贪图荣耀的心。因为荣耀对那当在各种谦卑中经受考验的人是不合法的。若贞洁之德是由天主所赐，\Quote{你为何夸耀，仿佛你没有领受呢？}如果你没有领受，\Quote{你有什么不是领受的呢？}但正因如此，显然那不是天主赐予你的——你不是唯独献给天主。让我们看看，属人的事物是否坚固真实。\par

% structure: p0043 source=h2 target=section reason=relative-heading-rank; base=h2

\section{第十四章 不蒙头给童贞女自身带来的危险}

有人报告说，当初这个问题刚被提出时，某人曾说：\Quote{我们如何邀请其他（\Emph{童贞女}）效仿同样的行为？}诚然，使我们快乐的是她们的数量，而非天主的恩宠和每个人各自的功德！在天主眼中，是\Emph{童贞女}装饰（或使）教会得称赞，还是教会装饰或称赞\Emph{童贞女}？（我们的反对者）于是承认\Quote{荣耀}是事情的根本。好，哪里有荣耀，哪里就有勾引；哪里有勾引，哪里就有强迫；哪里有强迫，哪里就有需要；哪里有需要，哪里就有软弱。所以，她们不蒙头，为的是因着荣耀被勾引，却又因着软弱导致的败坏而被强迫蒙肚皮。因为驱使她们的是争竞，而非宗教。有时那神——她们的肚皮——本身；因为弟兄们乐意承担对\Emph{童贞女}的供养。然而，不仅她们被败坏，她们还拖着\Quote{一条长绳的罪过。}因为，在被带到（教会）中间，因公开占有财产而得意，并被弟兄们以各种尊荣和慈善恩赐所充满之后——只要她们不跌倒——当任何罪过犯了时，她们就图谋一件与所享有的尊荣同样可耻的恶行。那就是：若蒙头是童贞女的公认标记，那么，若任何\Emph{童贞女}从童贞的恩宠中堕落，她为怕被发现，就永远不蒙头，并穿着那时确属他人的服饰行走。如今已成为无可置疑的\Emph{女人}，她们胆敢光着头走近天主。但那位\Quote{忌邪的天主和主}，祂说过：\Quote{没有掩盖的事不会被揭露的，}通常将这些女子带到众人眼前；因为除非被她们自己婴儿的哭声出卖，她们不会承认。但因她们\Quote{人数更多}，你岂不更当怀疑她们犯下更多罪吗？我要说（虽则我不愿说），一个害怕转变为\Emph{女人}的人，一旦秘密地转变了，还能在天主眼目下虚假地假装为\Emph{童贞女}，这是一件困难的事。再者，这样的人，为了怕被发现也是\Emph{母亲}，会对自己的子宫做出何等大胆的事！天主知道有多少婴儿，在长期被母亲抗拒之后，蒙祂帮助，经过怀胎而顺利完好地诞生！这样的\Emph{童贞女}总是最容易受孕，分娩最顺利，而且孩子确实最像他们的父亲！\par

这些罪行乃是强行与不情愿的\Emph{童贞}所招致的。不蒙头的贪欲本身就是不端庄的：它经验到某种非\Emph{童贞女}所该有的东西——即取悦人的追求，当然，而且是取悦\Emph{男人}。让她尽可尽量以诚实的心努力吧；她必因公开露面而陷于危险，当她被不可信赖的众多目光穿透，当她被指指点点所刺激，当她被过分疼爱，当她在不断的拥抱和亲吻中感到一股暖流。于是额头变硬；羞耻感消退；于是它松弛下来；于是学会了以另一种方式取悦的欲望！\par

% structure: p0046 source=h2 target=section reason=relative-heading-rank; base=h2

\section{第十五章 论迷惑}

不，真正而绝对且纯洁的\Emph{童贞}最怕的莫过于自身。甚至\Emph{女性的}目光它都回避相遇。它自身另有眼睛。它藉着头帕作为庇护，如同头盔，如同盾牌，以保护它的荣耀免受诱惑的打击，免受诽谤的堤坝，免受猜疑、私语、仿效，以及嫉妒本身。因为即使在异教徒中也有某种需警惕的东西，他们称之为着迷（或蛊惑），是过分赞美和荣耀的极其不幸的结果。我们有时在解释上将其归于魔鬼，因为仇恨善出自他；有时我们将其归于天主，因为对骄傲的审判出自祂，祂高举谦卑的人，贬抑得意的人。因此，越圣洁的\Emph{童贞女}，将害怕，即使以着迷之名，一方面害怕仇敌，另一方面害怕天主，前者的嫉妒性情，后者的监察之光，并将喜悦于只被自己和天主所知。但即使她被其他人认出，她也明智地堵住了诱惑的道路。因为谁有胆量用眼睛闯入被遮蔽的脸呢？一张没有感觉的脸？一张可谓阴郁的脸？任何邪恶的念头都会被这严厉本身打破。她，隐藏自己的\Emph{童贞}，由此甚至否认了自己的\Emph{女性身份}。\par

% structure: p0048 source=h2 target=section reason=relative-heading-rank; base=h2

\section{第十六章 特土良在表明其辩护与圣经、自然、纪律相合后，向童贞女们呼吁}

我们的意见的辩护即在于此，依照圣经、依照自然、依照纪律。圣经奠立法律；自然联合为之作证；纪律要求执行它。一个基于（单纯）意见的习惯，会为这三者中的哪一个辩护？或者相反观点的面目如何？圣经属于天主；自然属于天主；纪律属于天主。凡与此相悖的，不属于天主。如果圣经是不确定的，自然是显明的；关于自然的见证，圣经不能不确定。如果对自然有疑问，纪律指出更得天主认可的事。因为对祂而言，没有比谦卑更亲爱的，没有比端庄更悦纳的，没有比\Quote{荣耀}和讨人喜欢更可憎的。因此，愿你们视那些被天主认可的事为圣经、自然和纪律；正如你们奉命要\Quote{省察一切，并殷勤追随那更好的。}\par

同样，我们还需转向（童贞女）她们自己，以劝诱她们更乐意接受这些（建议）。我恳求你们，无论你们是母亲、姐妹还是童贞女儿——请允许我按你们年龄相称的称呼来称呼你们——请蒙上头帕：若是母亲，为你们的儿子之故；若是姐妹，为你们的弟兄之故；若是女儿，为你们的父亲之故。所有年龄都在你们身上面临危险。穿上端庄的全副武装；用羞怯的栅栏围绕自己；为你们的性别筑起壁垒，既不许自己的眼睛外出，也不许他人的眼睛进入。穿戴女性全套服饰，以保存童贞的身份。稍微掩饰你们内心的自觉，以便只向天主显明真理。然而，你装作新娘时并非掩饰自己。因为你已与基督成婚：你将肉身交托给祂；你将成熟许配给祂。行事要与你们新郎的旨意相合。基督正是那命令别人的未婚妻和妻子蒙头的那一位；当然，更命令祂自己的（新娘）。\par

% structure: p0051 source=h2 target=section reason=relative-heading-rank; base=h2

\section{第十七章 向已婚妇女的呼吁}

然而我们也劝诫你们，那些处于第二等（程度）谦虚的\Emph{女子}，你们已落入婚姻，切莫如此远离面纱的纪律，甚至在片刻之间，也不可因为你们不能\Emph{拒绝}它，就采取某种其他方式\Emph{使之无效}，既不遮盖也不裸露。因为有些人用头巾和羊毛带子，并不\Emph{遮盖}头部，而是将其束起；前面确实受到保护，但头部正常所在之处却裸露。另一些人则用亚麻小帽盖住大脑区域——我猜想是怕压迫头部——但并未完全覆盖到耳朵。如果她们的听力如此微弱，以至于不能透过覆盖物听见，我怜悯她们。让她们知道，整个头部构成\Quote{\Emph{女人}}。其界限和边界延伸到袍子开始之处。面纱的区域与散开时头发覆盖的空间相等；以便颈部也被环绕。因为正是\Emph{它们}必须被服从，为了它们，\Quote{权力}应当\Quote{在头上}：面纱是它们的轭。\Place{阿拉伯}的异教\Emph{女性}将是你们的审判官，她们不仅遮盖头部，还遮盖面部，如此彻底，以至于她们满足于一只眼睛自由，宁享受半光也不出卖整张脸。一个\Emph{女性}宁愿看也不愿被看。因此，某位\Place{罗马}女王说她们最不幸，因为她们更容易恋爱而非被爱；而她们却因免受第二种（且更常见的）不幸而更为\Emph{幸福}，即女性更容易被爱而非恋爱。异教学的谦虚确实更为单纯，且可以说是更为野蛮。\Emph{我们}的\Person{主}甚至通过启示，为面纱的延伸范围定下了尺寸。因为我们的一位姐妹曾受到天使这样的斥责，天使拍打她的脖子，仿佛在鼓掌：\Quote{优雅的脖子，理当裸露！你从头到腰揭开面纱固然很好，免得这脖子的自由对你无益！}当然，你对一人所说的，就是对所有人所说的。然而，那些在咏唱圣咏时，以及在任何提到天主之名时，仍不遮盖的人，将同样应得何等严厉的惩罚？（她们）甚至当即将花时间祈祷时，极其乐意地将一条流苏、一束绒毛、或任何线头放在头顶，并自以为遮盖了？她们如此错误地想象自己的头就这么小！另一些人，认为手掌显然比任何流苏或线头大，同样误用了她们的头部；如同某种（生物），与其说鸟不如说兽，虽有翅膀，头小，腿长，且直立行走。据说，当它需要躲藏时，它只把头——显然是整个头——伸进灌木丛，却暴露全身其余部分。因此，当它在\Emph{头部}安全，却在更大部位裸露时，它整个人连同头都被捉住。那些遮盖少于所需的人，也将遭遇同样的处境。\par

因此，必须在任何时候、任何地方，牢记法律行走，准备就绪，随时迎接每一次对天主的提及；祂若在心中，也将在\Emph{女性}的头上被认出。愿\Person{吾主耶稣基督}的平安与恩宠归给那些善意阅读这些（劝言）的人，归给那些更重视\OriginalTerm{效用}{Utility}而非\OriginalTerm{习俗}{Custom}的人；同样也归给\Person{塞普蒂米乌斯·特土良}，这篇论著是他的。\par

\SourceRecord{Translated by S. Thelwall.；Ante-Nicene Fathers；Vol. 4.；Edited by Alexander Roberts, James Donaldson, and A. Cleveland Coxe.；Buffalo, NY: Christian Literature Publishing Co.,；1885.；<http://www.newadvent.org/fathers/0403.htm>.}
